% Compile this single source twice with pdfLaTeX. All references and the
% integer verification of the numerical example are inline.
\documentclass[a4paper,11pt]{article}
\usepackage[T1]{fontenc}
\usepackage[utf8]{inputenc}
\usepackage{lmodern,microtype,amsmath,amssymb,amsthm,mathtools,booktabs}
\usepackage[left=25mm,right=25mm,top=23mm,bottom=24mm,headheight=14pt]{geometry}
\usepackage{xcolor,fancyhdr,fancyvrb}
\usepackage[hyphens]{url}
\usepackage[colorlinks=true,linkcolor=deep,citecolor=deep,urlcolor=deep]{hyperref}
\definecolor{deep}{HTML}{175A60}\definecolor{ink}{HTML}{182F3C}
\hypersetup{pdftitle={Counting decompositions in the original Erdos--Straus equation},pdfauthor={Keio and Shaitetsu}}
\pagestyle{fancy}\fancyhf{}\fancyhead[L]{\small\sffamily\color{ink} The Erd\H os--Straus equation at prime inputs}
\fancyhead[R]{\small\sffamily\color{ink} Keio and Shaitetsu}\fancyfoot[C]{\small\thepage}
\renewcommand{\headrulewidth}{0.3pt}
\newtheorem{theorem}{Theorem}[section]
\newtheorem{corollary}[theorem]{Corollary}
\newtheorem{lemma}[theorem]{Lemma}
\newtheorem{proposition}[theorem]{Proposition}
\setlength{\emergencystretch}{2.5em}\allowdisplaybreaks[1]\raggedbottom
\newcommand{\arxiv}[1]{\href{https://arxiv.org/abs/#1}{arXiv:#1}}
\begin{document}
\begin{center}
{\LARGE\bfseries\color{ink}Counting decompositions in the original\\[3pt] Erd\H os--Straus equation\par}
\vspace{6pt}
{\large Prime inputs, both divisibility types, and the exceptional set}\par\vspace{8pt}
Keio and Shaitetsu\quad\textperiodcentered\quad 20 November 2026
\end{center}
\begin{abstract}
For almost every prime $p$, the fraction $4/p$ has at least a constant times $(\log p)^2/\log\log p$ unordered decompositions of each of two types into three distinct unit fractions: exactly one denominator or exactly two denominators are divisible by $p$. All these decompositions can be sought with first denominator $x$ satisfying $(\log X)^8<4x-p\le(\log X)^9$, where $X<p\le2X$. For each fixed $A>0$, the exceptional set is $O_A(X/(\log X)^A)$, with the lower-bound constant allowed to depend on $A$. We obtain the bound by constructing divisors with prescribed residues and comparing the first two moments of their counts. It applies even to almost every prime in the thin set where the very first denominator fails. Scaling gives $\gg_A (\log n)^2/(\log\log n)^3$ solutions of each divisibility type for all but $O_A(X/(\log X)^A)$ integer inputs $X<n\le2X$. The universal existence assertion remains open.
\end{abstract}

\section{The quantifier in the original problem}
The Erd\H os--Straus conjecture asks whether, for every integer $n\ge2$, there are positive integers $x,y,z$ with
\begin{equation}\label{eq:original}
 \frac4n=\frac1x+\frac1y+\frac1z .
\end{equation}
It suffices to treat primes: if $n=ap$, multiply each denominator of a representation for $4/p$ by $a$. For $p=2$ use denominators $1,2,2$. For every odd prime $p$, a solution has one or two denominators divisible by $p$. Clearing denominators shows that at least one is divisible, and if all three were divisible then $4$ would be a sum of three unit fractions. The main theorem counts unordered solutions with three distinct denominators in each of the two cases. Continuity with the preceding prime-square question requires the input in \eqref{eq:original} to be $p$ itself.

Write $F_s(m,p)$ for the number of unordered three-term solutions for $m/p$ in which exactly $s$ denominators are divisible by $p$. For $m=4$, Elsholtz and Tao gave a lower bound $(\log p)^{(\log3)/2-o(1)}$ on a set of primes of relative density one for the total number of ordered solutions \cite[Theorem 1.8]{et}. They also bounded the mean over primes. Elsholtz and Planitzer distinguished large values along infinite subsequences of primes from density-one estimates for integer inputs \cite{ep}. The bounds here are for both types separately on the original prime inputs. A finite permutation factor does not affect comparison of exponents.

The existence and multiplicity problems have different exceptional sets. Vaughan's theorem bounds the number of integers up to $X$ with \emph{no} decomposition by $O(X\exp(-a(\log X)^{2/3}))$ for some $a>0$ \cite{vaughan,pw}. Our exceptional set concerns failure to have the stated \emph{number} of decompositions of each type in a specified range. A density-one multiplicity result leaves open the possibility of a prime with no decomposition anywhere. All logarithms are natural. The same argument works for every fixed numerator $m\ge2$.



\section{The divisor condition at a single prime}\label{sec:point}
For an odd prime $p$, any solution of \eqref{eq:original} has a smallest denominator $K$ in $p/4<K\le3p/4$. Write $c=4K-p$ and $N=pK$, so $(N,c)=1$. The remaining fraction $c/N$ splits if and only if
\begin{equation}\label{eq:exact}
 \exists\,d\mid N^2,\quad d>0,\quad d\equiv-N\pmod c .
\end{equation}
Indeed $(cy-N)(cz-N)=N^2$. A qualifying $d$ and its partner $N^2/d$ give $y=(N+d)/c$ and $z=(N+N^2/d)/c$. Each original solution supplies such a $K$. Conversely any split at such a $K$ solves \eqref{eq:original}; the three summands can be reordered. If $p>c$, both residual denominators necessarily exceed $K$.

\paragraph{What goes wrong with direct descent.}
At the scaled denominator $p^{a-1}K$ for $4/p^a$, the remainder is $c/(p^aK)$. Its signed divisors, obtained by dividing each divisor of $(p^aK)^2$ by $p^aK$, have residues
\[
 \left\{p^j\prod_{q\mid K}q^{e_q}\pmod c:
       |j|\le a,\ |e_q|\le v_q(K)\right\}.
\]
A split exists exactly when this set contains $-1$. The allowed exponent of $p$ shrinks when passing from $a=2$ to $a=1$. For instance $p=193$, $c=11$ and $K=51=3\cdot17$. Modulo 11, $p,17,3$ have exponents $-1,-1,-2$ in base 2. At $a=1$ the total exponent is between $-4$ and $4$, excluding $-1=2^5$. At $a=2$ the integer $p^2K$ is already $-1$ modulo 11. Thus the square input has a split at $193K$ while the original prime input has none at $K$. This computation concerns one candidate for $193$.

\paragraph{The prime from the square example.}
For $p=96\,616\,332\,289$, the least shift for $4/p^2$ was 163 in the preceding study. For $4/p$, the least shift is exactly 11. At $c=3$ one has $K=79\cdot305\,747\,887$, and every prime factor of $pK$ is $1\pmod3$. At $c=7$, $K=2\cdot11\cdot631\cdot1\,739\,957$. These four prime factors and $p$ are quadratic residues modulo 7. Both shifts fail by \eqref{eq:exact}. At $c=11$, set
\[
 K=24\,154\,083\,075,\qquad N=pK,\qquad d=183 .
\]
Then $d\mid K$ and $d\equiv-N\pmod{11}$; the two formulas after \eqref{eq:exact} give the decomposition of the original $4/p$. Direct integer checks for this example appear at the end.

\paragraph{A pointwise bound at successful denominators.}
Let $G=(\mathbb Z/c\mathbb Z)^\times$ have order $g$, with $(N,c)=1$ and $c\ge3$. Every residue which occurs among the divisors of $N^2$ occurs at least
\[
 (2g)^{1-g}\tau(N^2)
\]
times. To prove this, choose a divisor uniformly by choosing independently its exponent at each $q^e\parallel N$ from $0,\ldots,2e$. Every positive mass of a local residue is at least $1/(2g)$. In fact, if the order is $h$ and $2e+1<h$, the positive masses are $1/(2e+1)$; otherwise each is at least $\lfloor(2e+1)/h\rfloor/(2e+1)\ge 1/(2h)$. Convolve these laws successively. A factor contains the identity in its support, so the old support is contained in the new one. If the cardinality stays fixed, every translation from the new factor preserves the old support, and the least positive mass cannot decrease. If it grows, the least positive mass decreases by a factor of at most $2g$. There can be at most $g-1$ growth steps. Applying the bound to $-N$ in \eqref{eq:exact} proves the claim for successful shifts. Existence of a reachable residue remains the essential pointwise condition.


\clearpage\section{A lower bound for the number of decompositions at prime inputs}
\label{sec:mult}
Fix an integer numerator $m\ge2$. For a prime $p\nmid m$ write $F_s(m,p)$
for the number of unordered positive representations of $m/p$ as three
unit fractions in which exactly $s$ denominators are divisible by $p$.
The following bound applies to each type, including on the difficult
progressions for $m=4$. Write $\kappa=2-2\log 2=0.61370\ldots$.

\begin{theorem}\label{thm:two-types}
For fixed $m$ and any fixed reduced arithmetic progression, all but
\[
 O\!\left(\frac{X}{(\log X)^{1+\kappa-o(1)}}\right)
\]
of its primes $X<p\le2X$ satisfy, simultaneously for $s=1,2$,
\[
 F_s(m,p)\gg_m\frac{(\log p)^2}{\log\log p}.
\]
The lower bound already holds for distinct-denominator solutions whose smallest denominator is
$K$ in the range $(\log X)^8<mK-p\le(\log X)^9$.
The exceptional-set constant and the $o(1)$ may depend on the progression
and on $m$.
\end{theorem}

\begin{proof}
Split the progression into reduced classes modulo its least common
multiple with $m$. Write $p\equiv b\pmod M$ for one such class,
$m\mid M$, and use the uniform probability measure on its primes in
$(X,2X]$. All estimates below are as $X\to\infty$ with $m,M$ fixed.
Put $L=\log\log X$ and introduce the ranges
\[
 Y_0=(\log X)^8,\quad Y_1=(\log X)^9,\quad
 Q=(\log X)^{40},\quad z=X^{1/1000}.\]
Write \(
 \mathcal P=\{q\text{ prime}:Q<q\le z\}\).
For $q\in\mathcal P$ set $w_q=1/(q-1)$. Mertens' theorem yields
\[
 \nu=\sum_{q\in\mathcal P}w_q=L-\log L+O(1),\qquad
 \sum_{q\in\mathcal P}w_q\log q\le(1+o(1))\log z.
\]
Let $\mathcal C$ initially contain all $Y_0<c\le Y_1$ such that
$c\equiv-b\pmod m$, so that $(c,m)=1$.

\emph{Equidistribution of a prime factor.}
For $c\in\mathcal C$ let
\[
 \Delta_c=\max_{a\in(\mathbb Z/c\mathbb Z)^\times}
 \left|\sum_{q\in\mathcal P,\ q\equiv a\ (c)}w_q
           -\frac{\nu}{\varphi(c)}\right|.
\]
For every fixed $A>0$, $\sum_{c\in\mathcal C}\Delta_c\ll_A L^{-A}$.
Indeed the prime-counting form of Bombieri--Vinogradov \cite[Theorem 7.1]{harper} gives, uniformly for
$t\ge Q$, an aggregate error $O_H(t/(\log t)^H)$ for the prime counts
in reduced classes at moduli $c\le Y_1$: here
$Y_1\le t^{1/2}/(\log t)^{B(H)}$ for large $X$.
Partial summation with weight $1/(t-1)$ costs
$O_H((\log Q)^{1-H})$, including both endpoints. Replacing the common
logarithmic-integral main term by $\nu$ uses the prime number theorem
with an arbitrary logarithmic saving and
$\sum_{c\le Y_1}1/\varphi(c)\ll\log(2Y_1)$.
Retain the moduli $\mathcal C_0$ with $\Delta_c\le1/\varphi(c)$.
The discarded total reciprocal totient is at most $\sum\Delta_c$.
Consequently
\begin{equation}\label{eq:goodc}
 H=\sum_{c\in\mathcal C_0}\frac1{\varphi(c)}\gg_m L .
\end{equation}
For this lower bound it suffices to sum $1/c$ over the original
progression. At any good $c$ the $w$-mass of each prescribed unit
class lies between $(\nu-1)/\varphi(c)$ and $(\nu+1)/\varphi(c)$.

\emph{Two constructions.}
For a shift $c\in\mathcal C_0$, set $K=(p+c)/m$, $N=pK$. We have
$(N,c)=1$ and $p>c$. A positive divisor $e\mid N^2$ congruent to
$-N\pmod c$ gives the two remaining denominators $(N+e)/c$ and
$(N+N^2/e)/c$. Both exceed $K$.
Choose a set $S\subseteq\mathcal P$ of $t$ distinct primes dividing $K$
and subject to $d_S=\prod_{q\in S}q\le X^{1/10}$.
We use each of the following rules separately.
\begin{itemize}
\item For type $2$, choose exponents $a_q\in\{-1,1\}$ such that
$R=\prod_{q\in S}q^{a_q}\equiv-1\pmod c$, and put $e=NR$.
The two paired divisors each have $p$-valuation one.
\item For type $1$, choose exponents $a_q\in\{1,2\}$, at least one
equal to $1$, such that $R=\prod_{q\in S}q^{a_q}\equiv-m^{-1}\pmod c$.
Take $e=K^2/R$. It satisfies $e\equiv-mK^2\equiv-pK\pmod c$.
The paired divisors have $p$-valuations zero and two; exactly one
of the resulting denominators is divisible by $p$.
\end{itemize}
Both constructions give integral paired divisors because $q\mid K$ on $S$.
For type $2$, inversion of every exponent merely interchanges the
residual denominators, accounting for at most a factor two. The
type $1$ divisor of $p$-valuation zero is uniquely determined from the
unordered pair. Unique factorization recovers $S$ and its exponents
in either rule, up to the indicated interchange. Different shifts
have different smallest denominators. All denominators are distinct:
$e=N$ would force $c\mid2$, whereas $c>2$.

\emph{First moment.}
Write $W_{s,t}(p)$ for the total number of tuples supplied by rule $s$
before quotienting by interchange, summed over $c\in\mathcal C_0$.
We use every integer $|t-2\nu|\le\sqrt\nu$.
For one $c,s,t$ let $\mu_c$ be the sum over all tuples obeying the
residue and product conditions, with the requirement $d_S\mid p+c$
replaced by the weight $w(S)=\prod_{q\in S}w_q$. Uniformly,
\begin{equation}\label{eq:mulower}
 \mu_c\gg\frac{(2\nu)^t}{t!\,\varphi(c)} .
\end{equation}
To prove this, order the primes temporarily and fix an allowed exponent
pattern. Reserve one position whose exponent is $1$ or $-1$; type $1$
has at least one by construction. Select the other $t-1$ primes
independently with probabilities $w_q/\nu$. Their expected sum of
logarithms is $(2+o(1))\log z$. By Markov's inequality, a proportion
at least $9/10$ has product at most $X^{1/10}/z$ for large $X$.
The chance that any two selected primes coincide is at most
$O(t^2/((Q-1)\nu))=o(1)$. Whatever these choices, the remaining
prime must occupy exactly one reduced class modulo $c$. Its available
mass is at least $(\nu-1)/\varphi(c)-(t-1)/(Q-1)\gg\nu/\varphi(c)$.
Its size is at most $z$, giving the required product bound. There are
$2^t$ exponent patterns for type $2$ and $2^t-1$ for type $1$.
Division by $t!$ proves \eqref{eq:mulower}.

\emph{Replacement of prime counts by local weights.}
For either rule put $U=\sum_{c\in\mathcal C_0}\mu_c$. For the fixed
$t$ under consideration, the prime counting measure gives
$\mathbb E W_{s,t}=U+O_A((\log X)^{-A})$ for any fixed $A$.
In the second moment consider two tuples. With distinct shifts $c,c'$
a common $q$ would divide $c-c'$, impossible since $q>Q>Y_1$.
Otherwise, the congruences at the primes of $d=d_Sd_{S'}$ specify
one reduced class modulo $d$. At the same shift,
$d=\operatorname{lcm}(d_S,d_{S'})$ and the same conclusion holds.
In both cases the compatibility with $p\equiv b\pmod M$ follows from
$(d,M)=1$. The local weight for the combined event equals
$1/\varphi(d)$. The error after all replacements is again
$O_A((\log X)^{-A})$ in normalized expectation. In detail,
$d\le X^{1/5}$ by the product cutoff. For fixed $d$, there are
at most $O(Y_1^2\,3^{2t}4^t)$ ordered pairs of tuples: choose
the shifts, assign each prime to the first set, the second, or both,
and choose its exponents. This multiplicity is $(\log X)^{O(1)}$.
Apply Bombieri--Vinogradov to the moduli $Md$, retaining its maximum
over reduced classes, and choose its logarithmic saving arbitrarily
large before dividing by the number of primes. The prime number
theorem in the fixed progression replaces its count by
$(\operatorname{Li}(2X)-\operatorname{Li}(X))/\varphi(M)$ with any
power of logarithmic saving. Only polylogarithmic total model weight
is involved, as follows already by dropping every residue condition
and using $t=O(L)$ in the moment bounds below.

\emph{Variance.}
At two distinct shifts the model second moment is at most
$\mu_c\mu_{c'}$: overlapping supports are excluded and all remaining
weights factor. This bound uses the exact model means with their
cutoffs. At one shift let $r=|S\cap S'|$. For $r<t$, discard the
product cutoff and upper-bound the sum over distinct primes by one
allowing repetitions in its $2t-r$ positions. Divide the ordered
sum by $r!(t-r)!^2$, accounting for the three membership groups.
For each fixed pair of exponent patterns keep one private prime
on either side free. Each private congruence imposes at most
$R(c)=4\cdot2^{\omega(c)}$ reduced classes: the exponent is $1,-1$,
or $2$, and the kernel of squaring modulo $c$ has at most $R(c)$
elements by the Chinese remainder theorem. Thus the diagonal
contribution for $r<t$ is at most an absolute constant times
\[
 \frac{R(c)^2 4^t}{\varphi(c)^2}
 \frac{\nu^{2t-r}}{r!(t-r)!^2}.
\]
For $r=t$ discard the second condition. A first-pattern position of
exponent $1$ or $-1$ always exists; retaining it gives at most
$O(4^t\nu^t/(t!\varphi(c)))$. These upper bounds apply after
replacement by weights, so freeing primes introduces no primality
correlation.

Since $U\gg H(2\nu)^t/t!$, summing over $c$ and using
\[
 \sum_{r<t}\frac{(t!/(t-r)!)^2}{r!\,\nu^r}\le e^{t^2/\nu}
\]
bounds the relative variance by
\[
 \frac{\operatorname{Var}W_{s,t}}{U^2}\ll
  e^{t^2/\nu}\frac{1}{H}
   \max_{c\in\mathcal C_0}\frac{R(c)^2}{\varphi(c)}
  +\frac{t!}{H\nu^t}+O((\log X)^{-10}).
\]
The elementary estimates $2^{\omega(c)}=c^{o(1)}$ and
$\varphi(c)=c^{1-o(1)}$, together with $c>Y_0$, bound the first
term by $\exp((-4+o(1))L)$. By Stirling's formula the second term
is at most $\exp((-2+2\log2+o(1))L)$ throughout our window of $t$.
Chebyshev's inequality and a union bound over the two rules and
$O(\sqrt L)$ support sizes now give the exceptional-set estimate in
the theorem. Outside that set, each $W_{s,t}$ exceeds $U/2$. Supports
of different cardinalities give different representations, so we
may sum their counts. Stirling's formula in a central window gives
\[
 \sum_{|t-2\nu|\le\sqrt\nu}\frac{(2\nu)^t}{t!}\gg e^{2\nu}.
\]
Together with \eqref{eq:goodc} and $e^{2\nu}\asymp
(\log X)^2/L^2$, this gives $F_s(m,p)\gg_m(\log X)^2/L$.
\end{proof}


\subsection{An arbitrary logarithmic saving in the exceptional set}
The large prime factors used in the count permit the use of several disjoint intervals of shifts. This gives a stronger exceptional-set estimate for the multiplicity bound.
\begin{corollary}\label{cor:multamps}
Fix $A>0$, $m$ and the reduced prime progression. Outside a set of
$O_{A,m,M}(X/(\log X)^A)$ primes in $(X,2X]$, both counts in Theorem~\ref{thm:two-types} are $\gg_{A,m}(\log p)^2/\log\log p$ in the same first-denominator range.
\end{corollary}
\begin{proof}
Fix an integer $k$ with $1+k\kappa>A$. Subdivide the
shift interval at $(\log X)^{8+j/k}$, $0\le j\le k$.
Use $Q=(\log X)^{1000k}$, $z=X^{1/(1000k)}$, and support
cutoff $d_S\le X^{1/(10k)}$. The mean $\nu$ remains
$L-\log L+O_k(1)$. Each of the $k$ intervals has reciprocal shift
mass $\gg_m L/k$ by the same argument as \eqref{eq:goodc}.
The product cutoff in the first moment still holds by Markov's inequality,
since the expected sum is $(2+o(1))\log z$. For a fixed type and $t$ in
the central window, let $W_j$ count its tuples in interval $j$, and
let $U_j$ denote its exact model mean. Write $M_{j,a}$ for the exact
model moment of order $a=0,1,2$. Thus $M_{j,0}=1$, $M_{j,1}=U_j$,
and $M_{j,2}$ sums compatible pairs of local configurations with weight
$1/\varphi(\operatorname{lcm}(d_S,d_{S'}))$. The proof of
Theorem~\ref{thm:two-types} gives
\[
 0\le M_{j,2}-U_j^2\ll
 U_j^2\exp((-\kappa+o(1))L).
\]
The left inequality follows also by interpreting the exact local
model as choosing independent uniform units modulo each $q\in\mathcal P$.

For each $a_1,\ldots,a_k\in\{0,1,2\}$ we have the approximate
factorization
\begin{equation}\label{eq:multamps}
 \mathbb E\prod_j W_j^{a_j}
  =\prod_j M_{j,a_j}+O_k((\log X)^{-100k}).
\end{equation}
Indeed at most $2k$ tuples occur, so the lcm of their support products
is $\le X^{1/5}$. Their prime congruences are independent in the
CRT sense if compatible. The Bombieri--Vinogradov replacement still
has arbitrarily large logarithmic saving: for a specified lcm there
are at most $Y_1^{2k}(2^{2k}4^{2k})^{2kt}$ configurations, a fixed
power of $\log X$. Thus it suffices to compare the joint exact model
moment with the product of the one-interval model moments.
Configurations on which no prime is shared between different
intervals contribute identically. A shared prime would divide the
nonzero difference between two shifts, of absolute value below
$Q$. Such a configuration is absent from the joint model. In the
product model, each shared $q$ pays the weight $w_q^2$ from any
two intervals in which it occurs. To bound the sum, discard support-size,
product and residue restrictions. For each set of at most two labelled
tuples within an interval, a present prime has at most eight states:
it occurs only in the first, only in the second, or in both, with
two, two and four exponent choices. Counting choices for all shifts
first, the total collision weight is at most
\[
 O_k\!\left(Y_1^{2k}
    \prod_{q\in\mathcal P}(1+8w_q)^k
    \sum_{q\in\mathcal P}w_q^2\right)
 \ll_k Y_1^{2k}e^{8k\nu}/(Q-1)
 \ll_k(\log X)^{-900k}.
\]
This proves \eqref{eq:multamps}.

Expand the product of the normalized centered squares and apply
\eqref{eq:multamps} term by term. Since $U_j\gg1$, the resulting
error is $O_k((\log X)^{-100k})$. The product of the model terms
is the product of their normalized variances. Hence
\[
 \mathbb E\prod_{j=1}^k\left(\frac{W_j-U_j}{U_j}\right)^2
 \ll_k \exp((-k\kappa+o(1))L).
\]
If every $W_j<U_j/2$, the product in this expectation exceeds
$4^{-k}$. For almost all primes at least one interval therefore
contributes $\gg_m (L/k)(2\nu)^t/t!$ for each $t$ and each type,
with the total exceptional fraction $\ll_k
\exp((-k\kappa+o(1))L)$ after summing over the $O(\sqrt L)$ choices.
Different support sizes give different representations even when their
productive intervals differ. Summing in $t$ as before proves the result.
\end{proof}


\subsection{Returning to the integer input}
\begin{corollary}\label{cor:integer}
Fix $A>0$ and $m\ge2$. For all but $O_{A,m}(X/(\log X)^A)$ integers $n\in(X,2X]$, the fraction $m/n$ admits
\[
 \gg_{A,m}\frac{(\log n)^2}{(\log\log n)^3}
\]
distinct unordered representations of each of the following two kinds. In the first, exactly one of the three denominators is a multiple of $n$; in the second, exactly two are multiples of $n$. All three denominators in the constructed representations are distinct.
\end{corollary}
\begin{proof}
Put $L=\log\log X$ and $Y=X^{1/L}$. Apply Corollary~\ref{cor:multamps} with a sufficiently large fixed saving exponent $B>A+4$, and call a prime good if both global counts exceed the resulting fixed constant times $(\log p)^2/\log\log p$ with distinct denominators. Summing dyadically and by partial summation, the bad primes satisfy
\[
 \sum_{Y<p\le2X,\ p\ {\rm bad}}1/p\ll_{B,m}(\log Y)^{1-B}.
\]
At most $2X$ times this quantity of the integers in question have a bad prime divisor greater than $Y$.

Those having no prime divisor greater than $Y$ are also few. With $D=A+4$ and $\sigma=1-D/\log Y$, their number is at most
\[
 (2X)^\sigma\prod_{q\le Y}(1-q^{-\sigma})^{-1}
 \ll_D X e^{-D L}\log Y .
\]
Indeed $\sum_{q\le Y}(q^{D/\log Y}-1)/q\ll_D (\log Y)^{-1}\sum_{q\le Y}(\log q)/q=O_D(1)$, and the higher prime powers in the logarithm of the product contribute $O(1)$ for $\sigma\ge3/4$. Both exceptional counts are $O_{A,m}(X/(\log X)^A)$. Every remaining $n$ has a good divisor $p>Y$. Multiply the three denominators of each of its solutions by $n/p$. The map is injective for this $p$, preserves strict ordering, and $n\mid(n/p)d$ holds exactly when $p\mid d$. The stated bound follows from $\log p\ge(\log X)/L$.
\end{proof}


\section{Primes for which the first denominator fails}\label{sec:thin}
Let $\mathcal E$ denote the set of primes $p\equiv1\pmod{12}$ such that every prime factor of $(p+3)/4$ is $1\pmod3$. For an odd prime $p>3$ this is exactly the event that the smallest possible denominator $\lceil p/4\rceil$ fails. At $p\equiv3\pmod4$ the remainder is a single unit fraction, which splits. At $p\equiv1\pmod4$ the shift is $3$, $3\nmid pK$, and a signed divisor is $-1$ modulo $3$ precisely when some prime factor of $pK$ is $2\pmod3$. Its absence forces $p\equiv1\pmod3$ and $K$ odd. Thus every $p\in\mathcal E$ is also $1\pmod{24}$.

The counting function of $\mathcal E$ has order $X/(\log X)^{3/2}$. For the lower bound take the primitive form $F(u,v)=u^2+uv+v^2$ in the theorem of Fuchs, Hsu, Rickards, Schindler and Stange \cite[Theorem 1.1]{forms}, with shift $-3$ and multiplier $4$. Its discriminant is $-3$, its leading coefficient is one and $(-3,4)=1$. Primitively represented primes $p=4F(u,v)-3>3$ have $F$ odd and prime to $3$. Any odd prime $q\equiv2\pmod3$ dividing $F$ would force $q\mid u,v$ or a nontrivial cube root of unity modulo $q$. Thus these primes belong to $\mathcal E$. The upper bound follows by sieving the class $0$ at small primes and also $-3$ at primes $2\pmod3$, omitting $2,3$; the sieve is detailed in Section~\ref{sec:windowupper}.

Since $2-2\log 2>1/2$, the exceptional count of Theorem~\ref{thm:two-types}, summed over dyadic intervals, is $o(X/(\log X)^{3/2})$. Consequently almost every prime in $\mathcal E$ still has the stated number of solutions of each type. The same argument applies to every subset whose counting function is $\gg X/(\log X)^{3/2}$. Using cumulative counts avoids requiring a lower bound on each individual dyadic interval.

\paragraph{The second candidate on the same bases.}
For almost every prime of $\mathcal E$ it is also the case that $K=(p+7)/4$ works, so that the least shift equals 7. Here is a direct bound for the exceptions. If some prime $q\equiv-1\pmod7$ divides $K$, the signed ratio $q$ provides a split at 7. Simultaneous failures of 3 and 7 therefore exclude $0$ at small primes, $-3$ at primes $2\pmod3$ and $-7$ at primes $-1\pmod7$. Apart from finitely many small primes these are distinct residues. The interval upper-bound sieve of Section~\ref{sec:windowupper}, now with at most three classes at each prime and a sufficiently small fixed positive power of $X$ as cutoff, gives $O(X/(\log X)^{1+1/2+1/6})$ exceptions. The exponent exceeds $3/2$.

This comparison explains why a large first shift on even powers need not persist at the prime itself. Fix a prime $\ell\ge7$, $\ell\equiv3\pmod4$. In the primitive-form theorem above prescribe $p=1$ modulo all primes from $5$ to $\ell-1$ and $p=-1$ modulo $\ell$. These are allowed progressions in \cite[Theorem 1.1(2)]{forms}: their modulus is coprime to 6, and both the prescribed class and its translate by 3 are units. This leaves $\gg_\ell X/(\log X)^{3/2}$ primes in $\mathcal E$. For $n=p^{2h}$, $h\ge1$, every earlier $c\equiv3\pmod4$, $0<c<\ell$, has $(p/c)=1$. Every odd prime $q\mid(n+c)/4$ has $(q/c)=(-c/q)=1$ by reciprocity, since $-c\equiv(p^h)^2\pmod q$ and $p>\ell$. If $2\mid(n+c)/4$, then $c\equiv7\pmod8$ and $(2/c)=1$. The sign $(-1/c)=-1$ forbids a split. At $\ell$, writing $N=n(n+\ell)/4$, the divisor $Np$ supplies a split. Thus the exact least shift is $\ell$ at \emph{every} even power. Removing the exceptional sets in the two preceding paragraphs leaves the same order of primes, at which the original $4/p$ has least shift 7 and at least $\gg(\log p)^2/\log\log p$ decompositions of each type. This is a simultaneous statement about the same base primes.


\section{An upper bound in the same denominator window}
\label{sec:windowupper}
Fix a numerator $m$ and a constant $B>0$, and write $H=(\log X)^B$.
For $X<p\le 2X$, let $V(p)$ count the unordered three-term
decompositions of $m/p$ whose smallest denominator $K$ has
$0<c=mK-p\le H$. We have
\[
 \sum_{X<p\le 2X}V(p)\ll_{m,B}X(\log X)^{1+o(1)},\qquad
 \sum_{\substack{X<p\le 2X\\p\in\mathcal E}}V(p)
     \ll_{m,B}X(\log X)^{1/2+o(1)}.
\]
Consequently $V(p)\le(\log p)^{2+o(1)}$ in relative density in both populations (for each fixed exponent tolerance). Together with the lower bound
in the polylogarithmic band from the second moment, this gives the
exact exponent $2$ for each divisibility type throughout any
polylogarithmic window containing that band. In particular, with $0<mK-p\le(\log p)^9$, the logarithm of the count of each type divided by $\log\log p$ tends to $2$ in relative density. The statement about $\mathcal E$ uses cumulative counts of its primes. The upper estimates also
hold for larger fixed $B$.

To prove the bounds note that $(c,mKp)=1$ once $X$ is large: a
common prime factor of $m,c$ would divide $p$, and the same holds
for a common prime factor of $K,c$. Normalize a divisor of $K^2$
by $K$ and reduce the resulting ratio to $u/v$. Then $(u,v)=1$
and $uv\mid K$. Thus $K=uvw$ with $u,v,w\ge1$. Conversely these
choices give a signed divisor $u/v$ available from $K$. After
possibly ordering the two residual denominators, Type II requires
$u/v\equiv-1\pmod c$, while Type I requires
$u/v\equiv-p\pmod c$. The latter simplifies because
$p=muvw-c$. For an upper bound we can therefore count all triples
satisfying one of the following rows, dropping the coprimality of
$u,v$:
\[
\begin{array}{c|c|c}
 & p=muvw-c &\text{constraint}\\ \hline
\mathrm{II} & K=uvw & u+v\equiv0\pmod c\\
\mathrm{I} & K=uvw & mv^2 w+1\equiv0\pmod c .
\end{array}
\]
All variables remain coprime to $c$. Each residual unordered
completion arises from such a triple. Different completions in a
fixed row require different triples, upon choosing one orientation
per pair.

In either row choose a largest variable $L$, and call the other two
$a,b$. Then $ab\le K^{2/3}\ll_m X^{2/3}$, and as $L$ varies,
$p=mabL-c$ is in one class modulo $mab$. If $L$ occurs in the
constraint, its possible residues modulo $c$ belong to at most
$2^{\omega(c)+1}$ classes (one for a linear congruence). For the
quadratic case this uses the bound on the number of square roots
of a unit modulo $c$, including the $2$-part, by the Chinese remainder
theorem. This upper bound is $(\log X)^{o(1)}$ uniformly for $c\le H$
by the elementary maximal-order bound for the divisor function.
For each possible class $p$ is restricted to an arithmetic
progression of modulus $d=mabc\ll X^{2/3+o(1)}$.

We use the following uniform sieve bound for a single class of $p$
modulo $d\ll X^{2/3+o(1)}$. In the full population the bound is
$\ll(X/d)(\log X)^{-1+o(1)}$; in $\mathcal E$ it is
$\ll(X/d)(\log X)^{-3/2+o(1)}$.
Indeed on the progression omit all primes $q\mid6d$, exclude
$p\equiv0\pmod q$ at every remaining $q\le z=X^{1/200}$,
and in $\mathcal E$ additionally $p\equiv-3\pmod q$ at
$q\equiv2\pmod3$. For an interval of progression indices of
length $Y\asymp X/d$ the large sieve for omitted classes gives
$\ll Y\prod_{q\le z}(1-\nu(q)/q)$ with these exclusions.
For example, in the usual denominator sum truncated at $\sqrt Y$,
the full Euler-product weights have mean $\log d'$ at most
$2\log z+O(1)$, less than $\frac12\log\sqrt Y$.
This preserves at least half the full product. Mertens' theorem in
progressions proves the asserted exponents. Removing primes dividing
$6d$ costs a factor $\ll(d/\varphi(d))^{O(1)}
\ll (\log\log X)^{O(1)}$. When $c=1$ all congruences modulo $c$
are read as vacuous.

Write $\lambda=1$ or $3/2$ according to the population. If $L$
occurs in the constraint, summing the progression bound over $a,b,c$
therefore gives at most
\[
 X(\log X)^{-\lambda+o(1)}
       \sum_{c\le H}\frac1c
       \sum_{a,b\ll_m X}\frac1{ab}
 \ll X(\log X)^{2-\lambda+o(1)}.
\]
If $L$ does not occur in the constraint, the sieve modulus is
only $d=mab$. We must bound the constrained sum of $1/(ab)$ over
$a,b,c$. If $\max(a,b)\ge2H$, fix $c$ and the smaller variable.
The larger variable is then confined to at most
$2^{\omega(c)+1}$ residue classes modulo $c$. In each class the
sum of its reciprocals above $2H$ and up to $O_m(X)$ is
$O((\log X)/c)$, by comparison with a harmonic series. The sum of
the smaller reciprocals is $O(\log X)$, giving the same final estimate.
If $a,b<2H$ then for each pair $a,b$, the shift $c$ divides the
fixed positive integer $u+v\le4H$ or $mv^2w+1\le 8mH^3+1$,
with the two participating variables being exactly $a,b$. The number
of its divisors is $(\log X)^{o(1)}$, uniformly. Thus this part
contributes only $(\log X)^{o(1)}(\sum_{a<2H}1/a)^2$ to the
constrained harmonic sum. The displayed mean bounds follow.

Lastly, Markov's inequality gives the density upper exponent using
the prime number theorem for the full population and
$\#\{p\le X:p\in\mathcal E\}\gg X/(\log X)^{3/2}$ for the thin
population. Dyadic summation is harmless, discarding $p\le\sqrt X$.

\paragraph{An explicit iterated logarithm.}
In fact the two mean bounds just obtained can be written as
$O_{m,B}(X(\log X)^{2-\lambda}(\log\log X)^4)$.
In the progression sieve, deletion of the prime factors of $6d$
costs $O((d/\varphi(d))^2)$: for $q\ge5$ the ratio
$(1-\nu(q)/q)^{-1}(1-1/q)^2$ is at most $1+O(q^{-2})$.
This accounts for at most $(\log\log X)^2$. Retain the root bound
under the harmonic sum over moduli; since
$2^{\omega(c)}=\sum_{j\mid c}\mu^2(j)$,
\[
 \sum_{c\le H}\frac{2^{\omega(c)+1}}c
 \ll (\log(2H))^2\ll_B(\log\log X)^2 .
\]
The small-pair contribution was only $(\log X)^{o(1)}$ to the
harmonic sum before the sieve factor, which is absorbed by
$(\log X)^2$. Thus, in the full and thin populations separately,
all but a relative $O(1/\log\log X)$ of the primes have
$V(p)\ll_{m,B}(\log p)^2(\log\log p)^5$. For dyadic intervals
the relative comparison for the thin population is made after summing
the exceptional counts up to $X$ and dividing by its cumulative
count $\gg X/(\log X)^{3/2}$.


\section{What remains for the original conjecture}
For each sufficiently large prime $p$, Section~\ref{sec:point} reduces the original problem exactly to finding at least one integer $K$ and divisor $d$ such that
\[
 p/4<K\le3p/4,\qquad d\mid(pK)^2,\quad d>0,\qquad
 d\equiv-pK\pmod{4K-p}.
\]
The moment estimates control a sum over primes. They do not supply the missing witness at an individual exceptional prime. For perspective, many candidates can fail at \emph{each} prime: when $p\equiv1\pmod4$, choose an even $u$ with $\sqrt{5p/4}<u<\sqrt{3p/2}$ and set $K=u^2-p$. Then $p/4<K<p/2$ and $K$ is odd. For every prime $q\mid K$, $p\equiv u^2\pmod q$ is a nonzero square, so for $c=4K-p$ one has $(q/c)=(-c/q)=(p/q)=1$. Also $(p/c)=(c/p)=(K/p)=1$ since $K\equiv u^2\pmod p$. Thus all divisors of $(pK)^2$ have positive symbol and the target $-pK$ has negative symbol. Distinct such $u$ supply $\gg\sqrt p$ certified failing candidates. This local test explains those failures, while \eqref{eq:exact} still has to succeed somewhere to prove the conjecture.

Corpus notes 042 and 050 \cite{vmrc} supplied the starting ideas of testing all divisors under a character and imposing several congruences compatibly. On the prime input the counting argument instead constructs many signed divisors and controls their coincidences. Standard external inputs are identified in the proofs. Current literature includes exact type-count identities \cite{jiang} and sieves for search depth \cite{dahan}. A comparison with the stated Elsholtz--Tao lower bound is possible because its hypotheses and conclusion are explicit. Historical priority of the estimates in this note remains unestablished.

\paragraph{Integer check for the numerical example.}
The following standalone Python 3 code checks every prime factor used for the two earlier failures, constructs the later decomposition, and verifies the original equation by integer multiplication.
\begin{Verbatim}[fontsize=\small]
from math import isqrt

def prime(n):
    return n >= 2 and all(n % j for j in range(2, isqrt(n) + 1))

p = 96616332289
assert all(prime(q) for q in (p,79,305747887,2,11,631,1739957))
assert (p+3)//4 == 79*305747887 and p%3 == 79%3 == 305747887%3 == 1
qs = [2, 11, 631, 1739957]
assert 4*2*11*631*1739957 == p+7
assert all(pow(q,3,7) == 1 for q in [p]+qs)
K, d = (p+11)//4, 183
N = p*K
assert N*N % d == 0 and 4*K-p == 11
e = N*N//d
assert (N+d)%11 == (N+e)%11 == 0
y, z = (N+d)//11, (N+e)//11
assert K<y<z and 4*K*y*z == p*(y*z+K*y+K*z)
\end{Verbatim}

\begin{thebibliography}{99}\small
\bibitem{et} C.~Elsholtz and T.~Tao, Counting the number of solutions to the Erd\H os--Straus equation on unit fractions, \emph{J. Aust. Math. Soc.} \textbf{94} (2013), 50--105. \href{https://doi.org/10.1017/S1446788712000468}{doi:10.1017/S1446788712000468}.
\bibitem{ep} C.~Elsholtz and S.~Planitzer, The number of solutions of the Erd\H os--Straus equation and sums of $k$ unit fractions, \arxiv{1805.02945}.
\bibitem{vaughan} R.~C.~Vaughan, On a problem of Erd\H os, Straus and Schinzel, \emph{Mathematika} \textbf{17} (1970), 193--198.
\bibitem{pw} C.~Pomerance and A.~Weingartner, Exceptions to the Erd\H os--Straus--Schinzel conjecture, \arxiv{2511.16817}.
\bibitem{harper} A.~J.~Harper, \emph{Elementary analytic methods in number theory}, Lecture Notes 2, Section 7, especially Theorem 7.1 (Bombieri--Vinogradov). \href{https://warwick.ac.uk/fac/sci/maths/people/staff/harper/elementarymethodsnotes2.pdf}{Lecture notes}.
\bibitem{forms} E.~Fuchs, C.~Hsu, J.~Rickards, D.~Schindler and K.~E.~Stange, Primes represented by shifted quadratic forms: on primitivity and congruence classes, \emph{Acta Arith.} \textbf{222} (2026), 371--391. \href{https://doi.org/10.4064/aa250502-24-10}{doi:10.4064/aa250502-24-10}.

\bibitem{jiang} X.~Jiang, Counting Type I and Type II solutions of the Erd\H os--Straus equation: concentration, local obstructions, and a pointwise comparison, preprint (2026), \arxiv{2609.09204}.
\bibitem{dahan} B.~Dahan, Sieve dimension and search depth for the Erd\H os--Straus conjecture, $n\equiv1\pmod{24}$, preprint (2026), \arxiv{2608.24035}.
\bibitem{vmrc} QEEDDD, \emph{Verifiable Mathematical Reasoning Corpus}, nos.~042 and 050.
\href{https://github.com/QEEDDD/Verifiable-Mathematical-Reasoning-Corpus/tree/c2ac34ea4622778187173f41bb1996deee8081b7}{Repository snapshot \texttt{c2ac34e}}.
\end{thebibliography}
\end{document}
